% =====================================================================
%  5. Antecedentes
% =====================================================================
\section{Antecedentes}
\label{sec:antecedentes}

Esta sección traza la trayectoria de los trabajos que preceden al problema, con un
criterio: interesa qué se intentó, qué quedó resuelto y dónde se detuvo cada línea.
El estado actual detallado de cada una se desarrolla en la sección~\ref{sec:estado-arte}.

\subsection{Del instrumento de revisión al aseguramiento automatizado}
\label{sec:ant-calidad}

La preocupación por la calidad de los recursos educativos digitales produjo primero
\textbf{instrumentos de revisión humana}. El
\emph{Learning Object Review Instrument} \parencite{leacock2007framework} fijó nueve
dimensiones evaluables por revisores formados, con un compromiso explícito entre
validez y eficiencia: se diseñó para que un evaluador produjera una reseña útil en
pocos minutos, agregable con la de otros bajo un modelo de participación convergente.

La segunda etapa trasladó la pregunta del recurso individual al
\textbf{repositorio}. \textcite{clements2015open} revisaron 82 artículos publicados
entre 1994 y 2014 y construyeron un marco (LORQAF) que clasifica los enfoques de
calidad según sean de política, tecnológicos o sociales. Su inventario de quince
instrumentos específicos delimita con precisión hasta dónde llegaba la automatización
en ese momento: pruebas automáticas de metadatos, verificación de enlaces, sistemas
de recomendación y marcado colaborativo de problemas. Ningún instrumento verificaba
el contenido.

La tercera etapa intentó \textbf{integrar} las vías automática y humana.
\textcite{tabares2017modelo} propusieron un modelo por capas --- gestión, revisión de
expertos y percepción de usuarios --- aplicado a 46 objetos de la Federación de
Repositorios de Objetos de Aprendizaje de Colombia. Su capa automática incorpora una
métrica de coherencia basada en reglas y similitud de coseno entre campos de texto
libre, con umbral de 0,7. Es el antecedente más cercano al uso de similitud semántica
como señal de calidad en un contexto educativo, y también el que muestra su límite:
la coherencia que mide es entre \emph{metadatos}, no entre los contenidos de
documentos distintos.

\begin{estado}
\small
\textbf{Dónde se detuvo esta línea.} Tres etapas, un mismo límite: lo automatizado
cubre la forma (metadatos, enlaces, disponibilidad) y lo sustantivo queda en manos de
revisores humanos. El problema de la consistencia \emph{entre} documentos de una
misma colección no se aborda en ninguna de las tres.
\end{estado}

\subsection{De la verificación de hechos a la inferencia documental}
\label{sec:ant-verificacion}

En paralelo, el procesamiento de lenguaje natural construyó la maquinaria para
decidir si una afirmación está respaldada por un texto. \textcite{thorne2018fever}
formalizaron la tarea con FEVER, separando tres clases --- respaldado, refutado, sin
información suficiente --- y adoptando un criterio de anotación conservador: sin
evidencia explícita, la afirmación se etiqueta como \enquote{sin información
suficiente} aunque el anotador crea saber la respuesta. Su mejor sistema alcanzó
31,87\,\% de exactitud exigiendo que la evidencia recuperada fuera la correcta, frente
a 50,91\,\% sin exigirla: el sistema acertaba la etiqueta con evidencia equivocada en
una fracción considerable de los casos.

Un hallazgo lateral de ese trabajo anticipa el problema de esta investigación: los
autores encontraron \textbf{información inconsistente entre páginas de la propia
fuente de referencia} --- Pakistán descrito con el 41.\textsuperscript{o} y el
42.\textsuperscript{o} PIB en dos páginas distintas
\parencite[p.~8]{thorne2018fever}. La inconsistencia intracorpus aparece allí como
anomalía observada; aquí es el objeto de estudio.

\textcite{koreeda2021contractnli} llevaron la tarea al nivel documental con
contratos: hipótesis fijas evaluadas contra documentos largos, con identificación de
la evidencia que las sustenta. Su resultado central para este proyecto es de
dificultad: con BERT-large, el $F_1$ de la clase contradicción fue de 0,357 frente a
0,834 en la clase de implicación. Con evidencia perfecta provista por un oráculo, ese
$F_1$ sube a 0,62--0,66, lo que sitúa la recuperación de evidencia como el cuello de
botella de la tarea.

\subsection{De la recuperación léxica a la generación anclada en evidencia}
\label{sec:ant-rag}

La línea de recuperación de información llega a este proyecto por una vía indirecta.
El marco probabilístico de relevancia y su función BM25 \parencite{robertson2009probabilistic}
fijaron durante dos décadas el punto de comparación para cualquier sistema de
recuperación textual. La recuperación densa de \textcite{karpukhin2020dense} mostró
que representaciones aprendidas podían superarlo por 9 a 19 puntos absolutos en
exactitud top-20, con una excepción instructiva: BM25 gana cuando la coincidencia
léxica entre consulta y pasaje es alta.

\textcite{lewis2020retrieval} combinaron recuperación y generación en una arquitectura
única, con dos propiedades que este proyecto aprovecha directamente: el conocimiento
se actualiza cambiando los documentos indexados sin reentrenar el modelo --- los
autores lo demuestran intercambiando índices de 2016 y 2018 --- y la memoria de texto
resultante es legible y editable por humanos. \textcite{gao2024retrieval} sistematizan
después la evolución del patrón y formulan la advertencia que define el riesgo de
aplicarlo aquí: la información errónea recuperada puede ser peor que la ausencia de
información.

\subsection{Del agente que razona al agente que debe ser evaluado}
\label{sec:ant-agentes}

\textcite{yao2023react} formalizaron la alternancia entre razonamiento y acción sobre
herramientas externas, con dos resultados pertinentes: en verificación de hechos las
trayectorias resultaron más ancladas en evidencia que las de razonamiento puro, y el
proceso es inspeccionable --- un humano puede distinguir el conocimiento interno del
externo y corregir el comportamiento editando los pasos de razonamiento. El mismo
trabajo documenta sus modos de fallo: errores de razonamiento en el 47\,\% de los
casos analizados, incluidos bucles, y búsquedas no informativas en el 23\,\%.

Dos líneas posteriores acotan el entusiasmo. \textcite{mohammadi2025evaluation}
argumentan que completar la tarea no basta como criterio y que, por ser los agentes
estocásticos, la misma tarea debe repetirse varias veces para estimar su fiabilidad.
\textcite{zhan2024injecagent} muestran que los agentes con herramientas son
vulnerables a instrucciones inyectadas en el contenido que leen: GPT-4 con el patrón
ReAct resultó vulnerable en el 24\,\% de los casos del escenario base y en el 47\,\%
con un \emph{prompt} de refuerzo del ataque. El dato más relevante para un sistema que
procesa documentos subidos por usuarios es que los campos con alta libertad de
contenido son los más vulnerables, porque la instrucción maliciosa se mezcla con el
texto normal.

\subsection{Antecedente interno del proyecto}
\label{sec:ant-interno}

El proyecto cuenta con dos antecedentes propios. El primero es el documento
\emph{EduCurator AI --- Curaduría asistida de contenidos educativos con supervisión
docente} (V0), que estableció el problema, el marco tecnológico y una primera
formulación metodológica, dejando marcados de forma explícita los puntos que
requerían bibliografía real. El segundo es el MVP funcional, descrito en la
Parte~IV, que implementa la arquitectura y permite ejecutar el flujo completo sobre
un corpus.

\begin{advertencia}
La existencia del MVP es un antecedente de \emph{desarrollo}, no de
\emph{investigación}. No aporta evidencia sobre el comportamiento del sistema: aporta
la condición material que hace posible recogerla.
\end{advertencia}
